\section{Introduction}

Retrieval-augmented generation (RAG) systems \citep{lewis2020rag} are increasingly deployed in high-stakes settings where factual accuracy is critical---medical question answering, legal research, and enterprise knowledge management among them \citep{nakano2021webgpt,gao2023enabling}. Yet these pipelines share a critical failure mode: when retrieved evidence is internally conflicting, topically irrelevant, or simply too sparse to ground a reliable answer, most systems still produce a fluent, confident-sounding response \citep{zhang2023siren,bang2023multitask}. Consider a user asking ``What is the recommended dosage of ibuprofen for children?'' when the retriever surfaces two passages that contradict each other---one citing 5--10 mg/kg and another citing 10--15 mg/kg. A standard RAG pipeline blends or arbitrarily selects from these passages and outputs a single authoritative-sounding figure, with no signal to the user that the underlying evidence is in conflict. This behavior represents a systematic miscalibration between the quality of retrieved evidence and the assertiveness of the generated response \citep{guo2017calibration,maynez2020faithfulness}.

Prior work on RAG robustness has addressed related but distinct problems. Retrieval quality studies \citep{retrievalrobust2024,nomiracl2024} measure how often systems retrieve relevant passages, but stop short of modeling the downstream response policy. Uncertainty estimation approaches \citep{kuhn2023semantic,kadavath2022language} estimate model confidence but operate on parametric knowledge rather than retrieval signals. Hallucination detection methods \citep{factscore2023,minicheck2024,li2023halueval} evaluate the fidelity of generated text post hoc, without altering the generation decision. Conflict-aware reading comprehension work \citep{chen2022rich,conflictqa2023,xie2023adaptive} focuses on extractive or structured tasks where the correct answer can be determined by span selection. No existing system treats the response policy itself---the choice of when to answer confidently, when to hedge, when to disclose conflict, and when to abstain---as the primary object of optimization. We address this gap with \textbf{\method{}}: a lightweight, training-free framework that uses feature-based conflict detection to select among a structured policy menu at inference time. \method{} does not require GPU training or API calls to large models; it operates on retrieval dispersion scores, lexical contradiction cues, and compact claim-support features computable with sentence-transformers and scikit-learn on a standard laptop.

This paper makes three concrete contributions:
\begin{itemize}
    \item \textbf{A policy-level evaluation framework} for RAG under conflicting and low-quality evidence, defining four response actions---forced answer, disclaimer, conflict disclosure, and abstention---and providing the first systematic comparison of these policies across diverse retrieval failure modes.
    \item \textbf{A feature-based conflict detection policy} that combines retrieval dispersion scoring, lexical contradiction detection, semantic claim support, and language-mismatch signals into a logistic decision layer, with an explicit abstention-accuracy Pareto analysis across threshold settings.
    \item \textbf{A reproducible evaluation suite} spanning four datasets (WhoQA, NoMIRACL nonrelevant subsets, ConflictQA, and FaithDial), 3,073 evaluated examples, and five baselines ranging from no-filter RAG to a GPT-4 verifier, with all prompts, feature extraction code, and result traces publicly released.
\end{itemize}
